\chapter{The Two Worlds}

The first constitutional act which gave the Workers' and Peasants' Government a territorial name was the Declaration of Rights of the Toiling and Exploited People issued by the third All-Russian Congress of Soviets in January 1918---the Bolshevik counterpart of the Declaration of the Rights of Man and of the Citizen promulgated by the French revolution. This proclaimed Russia to be a Republic of Soviets of Workers', Soldiers' and Peasants' Deputies, adding that ``the Russian Soviet Republic is founded on the basis of a free union of free nations, as a federation of Soviet national republics''. The form of words preserved the international intentions of the revolutionary regime. Revolution was essentially international; it implied the substitution of class war for war between rival Powers. But the promotion of world revolution was also a prime necessity of the struggling Soviet regime. It was the only weapon available to the Bolsheviks when confronted with the embattled imperialist Powers; and without revolution, at any rate in the principal belligerent countries, the regime could hardly hope to survive. Nor could any distinction be drawn between the two belligerent camps; both were equally standard-bearers of the capitalist order, which the revolution sought to destroy. Hence any conception of foreign policy other than propaganda for revolution was at first alien to Bolshevik thinking. Trotsky, the first People's Commissar for Foreign Affairs, observed epigrammatically: ``I will issue a few revolutionary proclamations to the peoples of the world, and then shut up shop.'' 
% 注：修复了 Republic 和 capitalist 的轻微 OCR 拼写错误 (Repubhc -> Republic, capitahst -> capitalist)。

External realities, however, soon dissipated this vision, and forced on the struggling Soviet republic the role of a national state in a world of national states. The appeal to the warring nations for peace negotiations had fallen on deaf ears. Something had to be done about relations with Germany, whose armies had penetrated deep into Russian territory, and were still conducting operations of war. One of the first acts of the new government was to conclude an armistice with the imperial German Government and to sue for peace. In February 1918 peace negotiations opened at Brest-Litovsk. Trotsky, who led the Soviet delegation, demonstratively abandoned the traditional practices of diplomacy, appealed to the belligerent peoples over the heads of their governments, openly carried on anti-war propaganda among the German troops, and embarrassed the German delegation by pressing the demand for ``peace without annexations or indemnities'', which Germany, in dealings with the western Allies, had purported to accept. 

But German intransigence and the overwhelming superiority of German arms presented an inescapable dilemma. Trotsky could not reconcile it with his revolutionary principles to sign a humiliating treaty with an imperialist Power---a course which Lenin came to regard as inevitable. On the other hand, his sense of reality did not allow him to support the demands of Bukharin and other ``Left communists'' for a renewal of ``revolutionary war''. He devised the formula ``No peace, no war''. When, however, the Germans, not impressed by this undiplomatic eccentricity, resumed their advance, the same dilemma recurred in still starker form. Trotsky reluctantly cast his vote with Lenin for the acceptance of what Lenin himself called a ``shameful'' peace, involving the abandonment of the Ukraine, and of other large areas of former Russian territory, and resigned his post as People's Commissar for Foreign Affairs. The treaty was signed on March 3, 1918, and the German advance stayed. Simultaneously with the Brest-Litovsk negotiations, informal---and fruitless---approaches were made to British, French and American representatives in the hope of soliciting western aid against the Germans. These overtures to capitalist governments, no less than the signature of the Brest-Litovsk treaty, were bitterly resented, as a derogation from the international principles of the revolution, by a substantial minority of the party central committee, headed by Bukharin; and all Lenin's influence was required to secure approval for them. 
% 注：清理了 "/||of the Ukraine", "Iterritory", "(Foreign Affairs", "11 and the ' Brest-Litovsk" 等乱码。

The lessons of military impotence had now been borne in on the Bolshevik leaders. On February 23, 1918, even before the Brest-Litovsk treaty was signed, the Red Army, originally called ``the Workers' and Peasants' Red Army'' came into being; the date has since been celebrated annually as the birthday of the Red Army. Its name was intended to indicate its international revolutionary character and purpose. But the proclamation announcing its foundation was headed ``The Socialist Fatherland in Danger'', so that national as well as international consciousness presided over its birth. Trotsky was appointed People's Commissar for War with the task of organizing it. He was too much of a realist to suppose that an army could be built up out of raw untrained levies. His first response to the emergency was to recruit professional soldiers, former Tsarist officers, officially referred to as ``military specialists'', to train the new army. This expedient proved brilliantly successful. By the beginning of 1919 30,000 such officers had been enrolled. The Red Guard of 1917, which mustered barely 10,000 trained men, grew at the height of the civil war into a Red Army numbering five million. Trotsky himself displayed exceptional military talents. But he was also known for ruthlessness in his demand for unquestioning obedience and in punishing defaulters; and he had to extol virtues of military discipline which the revolution had set out to destroy. Desperate remedies were required in a desperate situation. 
% 注：去除了 " ' for unquestioning obedience " 中的杂乱单引号。

These expedients did not end the dangers besetting the regime, now transferred from Petrograd to its new capital, Moscow. Hostile ``White'' Russian military forces began to muster in different parts of the country. The German army remained, by agreement with a puppet national Ukrainian Government, in occupation of the Ukraine. The western governments, outraged by the revolution and by the Russian desertion of the Allies in the hour of their greatest need, decided to act. In March 1918 British, followed by French and American, forces occupied the northern port of Murmansk---ostensibly to protect the military stores accumulated there against a further German irruption. Meanwhile, the many thousand Czech prisoners of war in Russia, mainly deserters from the Austrian army, formed themselves into a Czech legion, and with the agreement of the Soviet Government set out for Vladivostok to embark there for the west. In Siberia the well-organized legionaries clashed with scattered and ineffective Soviet authorities, and---perhaps at first unwittingly---became a rallying-point for anti-Bolshevik forces. In April 1918 the Japanese Government, unwilling to be left out of the act, landed troops in Vladivostok, followed by British and American detachments two months later. In July British, French and American forces occupied Archangel. The survival of the Workers' and Peasants' Government in Moscow in the summer and autumn of 1918 seemed due not so much to its own strength as to the fact that the nations were linked in a life-and-death struggle on the western front, and had little thought for what happened elsewhere. 
% 注：去除了 "f ernments", "/ desertion", "I decided to act" 中混入的标记字母。

The collapse of Germany, and the armistice of November 11, 1918, gave a fresh turn to the screw. The incipient revolutionary situation in Berlin in the two months after the armistice, the successful revolutionary coups a few months later in Bavaria and in Hungary, as well as sporadic unrest in Britain, France and Italy, led the Bolshevik leaders to believe that the long-awaited European revolution was maturing. But events which offered hope and comfort to Moscow intensified the fear and hatred felt by the western governments for the revolutionary regime, and sharpened their determination to uproot it. The pretext that military operations in Russia were a subsidiary part of the war against Germany was perforce abandoned. Support was openly extended to Russian armies committed to the crusade against Bolshevism in Archangel, in Siberia, and in southern Russia. Now, however, a fresh complication occurred. The Allied troops, affected partly by war-weariness and partly by more or less outspoken sympathy for the workers' government in Moscow, were plainly unwilling to continue the fight. In April 1919 a mutiny in French naval vessels in Odessa forced the evacuation of the port. In Archangel and Murmansk the same end to the adventure was forestalled by the progressive withdrawal of the Allied troops. By the autumn of 1919 no Allied armed forces (except for Japanese and American contingents in Vladivostok) remained on Russian soil. 
% 注：将 "November II" 修正为正常的 "November 11"。

This set-back in no way modified the hostile intentions of the western Allies, who sought to compensate for the withdrawal of troops by an increased flow of military supplies, military missions, and verbal assurances to a number of would-be Russian ``governments'' arrayed against the Bolsheviks. The most promising of these was formed under the leadership of Kolchak, a former Tsarist admiral, who established some kind of authority over much of Siberia and began to move into European Russia; and in the summer of 1919 the Allied statesmen assembled in Paris for the peace conference entered into negotiations, which proved inconclusive, for the recognition of the Kolchak regime as the sole legitimate Russian Government. Denikin, a Tsarist general, enjoying strong Allied support, controlled southern Russia, overran the Ukraine, and in the autumn of 1919 reached a point 200 miles south of Moscow; and Yudenich, another general, mustered a White army in the Baltic for an assault on Petrograd. By this time, however, the Red Army had become an effective, though ill-equipped, fighting force. The various White armies were unable either to coordinate their efforts or to win the support of the populations in whose territories they operated. By the end of the year they were in headlong retreat. In January 1920 Kolchak was captured and executed by the Bolsheviks. By the spring of that year the White forces, except for a few isolated pockets of resistance, had been everywhere dispersed and destroyed. 
% 注：去除了段尾 "j the White forces" 和 "^ had been everywhere" 中的干扰符号。

The civil war hardened the stereotype, which had been shaping itself both in western and in Soviet thinking since October 1917, of two worlds confronting each other in irreconcilable contradiction---the capitalist world and the world of the revolution dedicated to overthrow it. After the collapse of German power in November 1918 Central Europe briefly became a bone of contention between the two worlds. The whiff of revolution in Berlin in January 1919 favoured the confident belief of the Bolsheviks that the death-knell of capitalism had sounded, and that the wave of revolution was in the process of spreading westward from Moscow. It was in this atmosphere that Lenin set out to realize an ambition nourished by him ever since the autumn of 1914---to replace the defunct Second or Social-Democratic International, which had split and destroyed itself on the outbreak of war through its abandonment of the principles of Marxism and internationalism, by a truly revolutionary Third or Communist International. It was the logical sequel of a decision taken by the party congress in March 1918 to replace the old party name Russian Social-Democratic Workers' Party, now sullied through its association with German social-democrats and Mensheviks, by the name Russian Communist Party (Bolsheviks). 
% 注：去除了 "^through its association" 中的标记。

Early in March 1919 more than 50 communists and sympathizers assembled in Moscow, of whom 35 held mandates from communist or near-communist parties or groups in 19 countries; many of these were small countries which had once formed part of the Russian Empire, and were now recognized as Soviet republics, including the Ukraine, Belorussia, the Baltic countries, Armenia and Georgia. The newly-founded German Communist Party sent a delegate with instructions to raise no objection of principle, but to seek a postponement of the creation of the International to a more propitious moment. Travel to Moscow from the west was virtually impossible. Groups in the United States, France, Switzerland, the Netherlands, Sweden and Hungary had given mandates to nationals resident in Moscow; the one British delegate had no mandate at all. The caution of the German delegate was overruled by the weight of enthusiasm. The arrival of a revolutionary Austrian delegate is said to have tipped the scale. The congress, constituting itself as the first congress of the Communist International (Comintern), voted a manifesto, drafted by Trotsky, tracing the decay of capitalism and the advance of communism since the Communist Manifesto of 1848; theses prepared by Lenin which denounced bourgeois democracy, proclaimed the dictatorship of the proletariat, and derided attempts to revive the discredited Second International; and finally a topical appeal to the workers of the world to bring pressure on their governments to end military intervention in Russia and recognize the Soviet regime. By way of providing the new-born International with an organization, the congress elected an executive committee (IKKI), and named Zinoviev as its president and Radek, now in a Berlin prison, as its secretary. A few days after the congress ended, the short-lived Hungarian Soviet republic was proclaimed in Budapest. 

The fact of the foundation of a Communist International was more important than anything done at its first congress. It was a dramatic announcement of the rift between the two worlds, and in particular of the rift which had declared itself within the international workers' movement. The founders of Comintern firmly believed that the workers of western countries who had lived through the fratricidal slaughter of the war---and especially the German workers, well-schooled in Marxism---would quickly abandon the national Social-Democratic and Labour Parties which had involved them in the holocaust, and rally to the cause of the international unity of the workers of the world proclaimed by Comintern. When this did not happen, and when the Second International even showed signs of revival, the set-back was attributed to corrupt and treacherous leaders who had betrayed their misguided followers. But the rift in western countries between a minority of committed communists and a majority of workers who remained faithful to ``reformist'' leaders perpetuated itself, and deepened as time went on. 
% 注：删除了 "4his did not happen" 中的排版错误，并清理了 "Tnrr iTITTif" 这种不可名状的 OCR 乱码区。

The breach was aggravated by unforeseen developments within Comintern itself. The outlook of its founders was genuinely international; they looked forward to the day when its headquarters might move to Berlin or Paris. But what happened in Moscow in March 1919 was not a fusion of national communist parties into one international organization, but the harnessing of a number of weak, embryonic foreign groups to an essentially Russian organization, whose resources and main motive force came necessarily and inevitably from the Russian party and the Soviet Government. Nor was this illogical. The promotion of international revolution had two aspects, which reinforced one another. It was an obligation of all Marxists, but it was also an important defensive weapon in the armoury of the hard-pressed Soviet regime. So long as the overthrow of capitalist domination elsewhere was seen as a condition of the survival of the revolutionary regime in Russia, there could be no incompatibility between the two elements; they were different facets of a single coherent and integrated purpose. But this meant that the commitment of foreign communist parties to Comintern had less strong foundations than the commitment which seemed obligatory in Moscow. 
% 注：修复了由于分栏原因造成的句子倒置断裂 ("Jnational communist parties into one international organization, The breach was aggravated by unforeseen developments within Comintern itself.")。

The remainder of 1919 was a period of civil war, Allied intervention and Soviet isolation. A brief respite followed the collapse of the White armies in the winter of 1919-1920; and it was in this interval, in April 1920, in preparation for the second congress of Comintern, that Lenin wrote his famous and influential pamphlet, The Infantile Disease of ``Leftism'' in Communism. The target of attack was a so-called Left opposition in communist parties which resisted ``compromises'' in the name of ``principles''; Lenin recalled, in particular, the opposition to Brest-Litovsk. Communists in western countries must participate actively in parliaments and trade unions, and not shrink from the compromises inherent in such participation. Mindful of hostile British intervention in the civil war, Lenin urged British communists to conclude ``electoral agreements'' with the Labour Party in order to ``help the Hendersons and the Snowdens to defeat Lloyd George and Churchill''. But this advice was tendered against a background of confidence in the early prospect of revolution. The tactical prescriptions of the pamphlet were deeply imbued with the need to enlighten the rank and file of workers' parties on the true character of their leaders, and to split the parties against the leaders. Henderson was to be supported ``as the rope supports the man who is being hanged''. It did not enter into Lenin's calculations that such tactics of compromise and manoeuvre might be continued, in default of international revolution, for years or decades. 
% 注：清理了段首的 "nil"、"'^11'intervention" 等标记。

At the end of April 1920 Pilsudski launched a Polish invasion of the Ukraine, occupying Kiev early in May; and the Soviet republic was plunged once more into a crisis as grave as that of the civil war. But this time the resistance was swifter and stronger. In June the Red Army counter-attacked. The defeat of the over-extended Polish forces became a rout, and at the beginning of August the Red Army entered Polish territory. These dramatic events coincided with the second congress of Comintern, which opened on July 19, 1920, with more than 200 delegates. These included, besides delegates of the small German Communist Party (KPD), delegates of the German Independent Social-Democratic Party (USPD), a war-time break-away from the German Social-Democratic Party (SPD), as well as of the French and Italian Socialist Parties; these three parties were divided among themselves on the question of adhesion to Comintern, and had come to the congress for enlightenment. Delegates also came from several British groups of the extreme Left, which decided to merge in a Communist Party of Great Britain (CPGB). The debates, against a background of victories of the Red Army, were full of confidence and excitement. The prescriptions of Lenin's pamphlet were not forgotten. Resolutions were passed urging communists to work in trade unions and bourgeois parliaments; and the British Communist Party was instructed---by a majority vote---to seek affiliation to the Labour Party. But the prevailing mood was now quite different. The congress appealed to the workers of the world not to tolerate ``any kind of help to White Poland, any kind of intervention against Soviet Russia''. World revolution was kept well in the picture. 
% 注：修复了多栏扫描造成的语序错乱 ("I that of the civil war. But this time the resi.stance was swifter of the Ukraine, occupying Kiev early in May; and the Soviet...") 以及清除了 "hnd stronger"、"/defeat"、"fat the beginning" 等噪点。

\begin{quote}
The Communist International [declared a manifesto of the congress] proclaims the cause of Soviet Russia as its own cause. The international proletariat will not sheath the sword till Soviet Russia becomes a link in a federation of Soviet republics of the world. 
\end{quote}
% 注：根据排版规范，将这段引言放入了 blockquote 环境，以增加层级感。

The ``21 conditions'' of admission to Comintern drawn up by the congress were designed to exclude waverers, and to make Comintern, not (like the Second International) a loose association of widely differing parties, but a single homogeneous, disciplined party of the international proletariat. Never had the prospect of world revolution seemed so bright and so near. 

While the congress debated, the Soviet leaders had to take a vital decision. Was the Red Army to stand on the Polish frontier, and offer terms of peace to Pilsudski? Or was it to continue its now almost unopposed advance on Warsaw and other industrial centres of Poland? Lenin declared for the advance, dazzled by the prospect that the Polish workers would welcome the Red Army as their deliverers from the capitalist yoke, and that revolution in Poland would open the gate-way to Germany and western Europe. Trotsky and Radek came out against him; Stalin seems to have shared their doubts, but was absent at the front at the time of the critical decision. Tukhachevsky, the brilliant commander who had led the counter-offensive, was all for the advance, and wished to make the Red Army the army of Comintern. Boldness and enthusiasm carried the day. By the middle of August the Red Army was deployed before Warsaw. Here, however, the major miscalculation in these proceedings was quickly revealed. The Polish workers did not stir; and Pilsudski successfully appealed for national resistance to the Russian invader. In the next few weeks, the Red Army suffered the humiliations of precipitate retreat which it had so recently inflicted on its adversaries. The armies finally came to rest at a point far to the east of the so-called ``Curzon line'', which had been recognized by the Allied governments, as well as by the Soviet Government, as the eastern frontier of Poland. Here an armistice was signed on October 12, 1920. The Soviet republic had paid a heavy price for its revolutionary optimism. 

The prestige of the Red Army was partially retrieved by the ease with which it repelled an attack by Wrangel, the last of the White generals, in southern Russia in the autumn of 1920. But the defeat in Poland had lasting repercussions on Soviet relations with the western world. The campaign had been based on a conviction that the Polish workers would revolt against their rulers, and in conjunction with the Russian forces instal a revolutionary government in Warsaw. The disappointment of that hope showed that the Polish workers, like those of western Europe, were still too deeply imbued with national loyalties to embrace the cause of international proletarian revolution. Elsewhere in Europe, while the workers continued to demonstrate sympathy and enthusiasm for the Russian revolution, they showed no alacrity to raise the banner of revolution in their own countries. In October the USPD decided by a narrow majority to merge with the KPD, leaving the rump of its membership, together with the largest German workers' party, the SPD, to nourish feelings of bitterness and resentment against the KPD and against Comintern. A little later the French Socialist Party transformed itself into the French Communist Party (PCF), leaving behind a substantial minority of dissidents; and a split in the Italian Socialist Party led to the creation of a small Italian Communist Party (PCI). These additions to the membership of Comintern were hailed as triumphs in Moscow. But they hardened the distrust of Comintern now prevailing in many sections of the workers' movement in the west. An attempted revolutionary coup in Germany in March 1921 (see p. 44 below) was a dismal failure. The post-war revolutionary wave in Europe was visibly receding. 
% 注：去除了 "I workers"、"l^istrust"、"||the workers'"、"jtary coup"、"'visibly receding" 中夹杂的乱码。

Another lesson could also be drawn from the military defeat in Poland. The Russian peasant who supplied the manpower of the Red Army, while stoutly defending the revolutionary cause in his homeland, had no stomach for a fight to transport revolution to other countries. The peasant, now beginning to revolt against the miseries and the devastation which were the aftermath of the civil war, was recalcitrant to hardships sustained in the name of international revolution. In the hard winter of 1920-21 peasant disturbances in central Russia concentrated the anxious attention of the leaders on domestic problems, and began insensibly to re-shape Soviet thinking about the western world. Visions of international revolution had been encouraged---almost imposed---by the traumatic experience of the civil war. Once this was overcome, the aim of international revolution, though not disavowed, was allowed quietly to recede into a more distant future. Security and stability were the paramount needs of the moment. In this mood, steps were taken, simultaneously with the introduction of NEP, to regularize Soviet relations with the non-Soviet world. 
% 注：此段恢复了严重的句子断裂 ("f)f the Red Army, while stoutly defending the revolutionary Another lesson could also be drawn from the military defeat in Poland. The Russian peasant who supplied the manpower cause in his homeland, had no stomach for a fight to transport U^Vs, revolution to other countries.") 并根据您的指示在相应部分加上了 \textbf{}。